重叠滑动窗口是基于持续同调的时间序列变化检测中常用的局部化机制。然而，同一原始观测被重复纳入相邻窗口会改变窗口特征的联合依赖结构，从而使全排列校准所要求的交换性条件可能失效。本文首先推导重叠窗口均值的精确协方差结构，并构造有限样本反例，说明原始序列独立同分布时，派生窗口序列仍可能不具交换性。随后，在固定窗口数与固定观测长度两种预算制度下，基于 28 个预先设定的零假设配置，系统比较朴素置换与依赖乘子自助法在拓扑表示和矩表示上的有限样本 I 类错误率。在名义显著性水平 0.05 下，朴素置换的逐格拒绝率最高达到 0.9996；依赖乘子法明显削弱了极端反保守性，但其逐格拒绝率仍分布于 0.0020--0.1412，显示出同时存在的保守区与反保守区。进一步地，形状变化与均值平移两类方向相反的备择使两种表示的尺寸标准化功效发生近乎完全的方向反转，表明检出能力具有明显的信号类型选择性。上述结果将重叠窗口下的推断有效性刻画为随重叠程度、原序列依赖、预算制度和表示而变化的有限样本校准表面，并说明表示敏感性与校准有效性需要分别评价，同时考察二者的有限样本交互。

\noindent\textbf{关键词}：持续同调；变点检测；重叠滑动窗口；交换性；有限样本校准。
